% Generated draft: requires mathematical and source review.
\documentclass[11pt]{article}
\usepackage[margin=25mm]{geometry}
\usepackage{fontspec}
\setmainfont{Latin Modern Roman}
\usepackage{amsmath,amssymb,mathtools,unicode-math}
\setmathfont{Latin Modern Math}
\usepackage{enumitem,needspace,xurl,hyperref}
\setlength{\parindent}{0pt}
\setlength{\parskip}{5pt}
\setlength{\emergencystretch}{3em}
\newcommand{\N}{\mathbb{N}}
\newcommand{\Z}{\mathbb{Z}}
\newcommand{\Q}{\mathbb{Q}}
\newcommand{\R}{\mathbb{R}}
\newcommand{\C}{\mathbb{C}}
\newcommand{\F}{\mathbb{F}}
\newcommand{\A}{\mathbb{A}}
\newcommand{\PP}{\mathbb{P}}
\newcommand{\E}{\mathbb{E}}
\newcommand{\eps}{\varepsilon}
\newcommand{\dd}{\,\mathrm{d}}
\begin{document}
\section*{UnsolvedMath ID: 30001948}
\textbf{Source problem code:} catalogue entry
\subsection{Definitions and mathematical statement}
A stratified pseudomanifold is a space decomposed into manifold strata, with local conical models whose cross-sections are called links. In the usual rational convention, such a space satisfies the Witt condition if every even-dimensional link \(L^{2k}\) of a singular stratum satisfies
\[
I^{\bar m}H_k(L;\Q)=0,
\]
where \(I^{\bar m}H_*\) denotes intersection homology with lower-middle perversity \(\bar m(c)=\lfloor(c-2)/2\rfloor\). The supplied record does not specify the coefficient convention or the precise category of stratified spaces; this definition explains the usual terminology rather than supplying omitted hypotheses.

For a closed smooth oriented manifold \(M\), the usual homology \(\widehat A\)-class is \(\widehat A(TM)\cap[M]\), where \(\widehat A(TM)\) is the rational characteristic class associated with the \(\widehat A\) multiplicative sequence. This smooth construction does not itself define the requested class on a singular space.

\textbf{Question.} For a Witt space \(X\), determine hypotheses under which a homology \(\widehat A\)-class can be defined topologically. When positive-scalar-curvature metrics on the links are used in its construction, determine its dependence on those metrics. Determine also whether a suitable spin Dirac operator defines a \(K\)-homology class whose Chern character is that homology \(\widehat A\)-class; symbolically, after making the required choices and constructions precise, ask whether
\[
\widehat A_*(X;g_{\mathrm{links}})
=\operatorname{ch}([D_X]).
\]
Here \(g_{\mathrm{links}}\) denotes the link metric data, and \([D_X]\) denotes the proposed class of a spin Dirac operator. These symbols express the requested constructions, not objects asserted to exist under the Witt condition alone.

The source leaves the necessary hypotheses as part of the question. In particular, it does not specify compactness, boundary conditions, orientation or spin data, regularity and compatibility of link metrics, the analytic domain of the Dirac operator, or the intended homology and \(K\)-homology theories. Consequently, the displayed identity is a conditional target requiring these specifications, rather than a fully determined assertion for every Witt space.
\subsection{Short English statement}
Determine when a Witt space admits a topologically defined homology \(\widehat A\)-class, how a construction using positive-scalar-curvature link metrics depends on those metrics, and whether the class is the Chern character of a spin Dirac \(K\)-homology class. The source does not supply the hypotheses needed to make all these constructions precise.
\subsection{Sources}
\begin{itemize}
\item Supplied local UnsolvedMath registry record: permanent ID \(30001948\), source problem code \texttt{catalogue entry}, title ``Topological \(\widehat A\)-Classes on Witt Spaces.'' Project registry locations are \texttt{data/unsolvedmath/problems.json} and \texttt{catalog/problems.json}; this draft uses the supplied complete record.
\item Supplied report citation: \emph{Stratified Spaces: Joining Analysis, Topology and Geometry} (2012), Oberwolfach Reports, pp.\ 3217--3286, DOI \texttt{10.4171/owr/2011/56}. Supplied source URL: \url{https://doi.org/10.4171/owr/2011/56}. The collection describes these reports as records of mathematical workshops and their open-problem discussions. No individual proposer is supplied; the proposed year is recorded as 2011.
\item Supplied literature source, labeled ``original OWR report'': \url{https://doi.org/10.4171/owr/2011/56}.
\item Literature assessment dated \texttt{2026-08-21}: the supplied curation labels the problem open and states that metric dependence and identification with a spin-Dirac \(K\)-homology Chern character remain unresolved in the cited stratified-space literature. This is supplied dated curation, not an independently verified literature assessment; time-sensitive status and cited claims require independent checking.
\item The registry labels the statement \texttt{corrected\_verified} and reports that its mathematical content was checked against the official report, with notation, context, and question form normalized. This is a supplied verification claim; the report has not been independently consulted for this draft.
\end{itemize}
\end{document}
